\section*{Hoofdstuk \ref{ch:g10:stats} --- Beschrijvende statistiek}

\begin{solution}{exo:g10:stats:1}
Som: $0+1+1+2+2+2+3+3+4+4+5+6+7+8+12 = 60$, dus is het \omterm{def:g10:stats:mean}{gemiddelde}
$\frac{60}{15} = 4$. De reeks staat al gesorteerd en $N = 15$ is oneven: de
\omterm{def:g10:stats:median}{mediaan} is de $8$ste waarde, $3$. \omterm{def:g10:stats:quartiles}{Variatiebreedte}: $12 - 0 = 12$.
\end{solution}

\begin{solution}{exo:g10:stats:2}
\omterm{def:g10:stats:series}{Frequenties}: $\frac{7}{50} = 0.14$, $\frac{9}{50} = 0.18$,
$\frac{8}{50} = 0.16$, $\frac{10}{50} = 0.20$, $\frac{6}{50} = 0.12$,
$\frac{10}{50} = 0.20$ (som $1$). \omterm{def:g10:stats:mean}{Gemiddelde} uitkomst:
\[
\bar x = \frac{7 \times 1 + 9 \times 2 + 8 \times 3 + 10 \times 4
+ 6 \times 5 + 10 \times 6}{50}
= \frac{7 + 18 + 24 + 40 + 30 + 60}{50}
= \frac{179}{50} = 3.58 .
\]
\end{solution}

\begin{solution}{exo:g10:stats:3}
$N = 12$ (even): de \omterm{def:g10:stats:median}{mediaan} is het \omterm{prop:g10:coordgeom:midpoint}{midden} van de $6$de en de $7$de waarde,
allebei gelijk aan $6$: \omterm{def:g10:stats:median}{mediaan} $6$.

$Q_1$: $\frac{12}{4} = 3$, dus is $Q_1$ de $3$de waarde: $4$.

$Q_3$: $\frac{3 \times 12}{4} = 9$, dus is $Q_3$ de $9$de waarde: $8$.

Boxplot: doos van $4$ tot $8$ met de medianenstreep in $6$, snorharen van
$1$ tot $11$.
\end{solution}

\begin{solution}{exo:g10:stats:4}
De $24$ resultaten tellen op tot $24 \times 11.5 = 276$. Met het nieuwe
resultaat is het totaal $276 + 19 = 295$ voor $25$ leerlingen: het nieuwe
\omterm{def:g10:stats:mean}{gemiddelde} is $\frac{295}{25} = 11.8$.
\end{solution}

\begin{solution}{exo:g10:stats:5}
Huidig totaal: $5 \times 12 = 60$. Een \omterm{def:g10:stats:mean}{gemiddelde} van $13$ over $6$ toetsen
vraagt een totaal van $78$, dus moet het $6$de resultaat $78 - 60 = 18$ zijn.
Een \omterm{def:g10:stats:mean}{gemiddelde} van $14$ zou $84 - 60 = 24 > 20$ vragen: niet haalbaar.
\end{solution}

\begin{solution}{exo:g10:stats:6}
\omterm{def:g10:stats:mean}{Gemiddelde} groter dan $20$: bijvoorbeeld $1,\ 2,\ 3,\ 10,\ 40,\ 50,\ 60$ ---
\omterm{def:g10:stats:median}{mediaan} $10$ (de $4$de gesorteerde waarde), \omterm{def:g10:stats:mean}{gemiddelde}
$\frac{166}{7} \approx 23.7$.

\omterm{def:g10:stats:mean}{Gemiddelde} kleiner dan $6$: bijvoorbeeld $0,\ 0,\ 0,\ 10,\ 10,\ 10,\ 10$ ---
\omterm{def:g10:stats:median}{mediaan} $10$, \omterm{def:g10:stats:mean}{gemiddelde} $\frac{40}{7} \approx 5.7$.

Het \omterm{def:g10:stats:mean}{gemiddelde} kan niet onder $\frac{40}{7}$ zakken: opdat de \omterm{def:g10:stats:median}{mediaan} (de
$4$de gesorteerde waarde) $10$ zou zijn, moeten de waarden op de plaatsen
$4$ tot $7$ alle $\geq 10$ zijn, zodat het totaal minstens $40$ is en het
\omterm{def:g10:stats:mean}{gemiddelde} minstens $\frac{40}{7}$ --- het tweede voorbeeld is extremaal.

Deze voorbeelden tonen dat het \omterm{def:g10:stats:mean}{gemiddelde} en de \omterm{def:g10:stats:median}{mediaan} grotendeels
onafhankelijk zijn: de ene kennen zegt weinig over de andere, en daarom
meldt een goede samenvatting beide.
\end{solution}

\begin{solution}{exo:g10:stats:7}
\emph{1.} \omterm{def:g10:stats:mean}{Gemiddelde}:
$\frac{9 \times 1800 + 12000}{10} = \frac{16200 + 12000}{10} = 2820$.
\omterm{def:g10:stats:median}{Mediaan}: gesorteerd zijn de $10$ lonen negen keer $1800$ en dan $12000$; de
\omterm{def:g10:stats:median}{mediaan} is het \omterm{prop:g10:coordgeom:midpoint}{midden} van de $5$de en de $6$de waarde, allebei $1800$:
\omterm{def:g10:stats:median}{mediaan} $1800$.

\emph{2.} De \omterm{def:g10:stats:median}{mediaan} ($1800$) beschrijft het typische loon: dat is wat $9$
werknemers op $10$ werkelijk verdienen. Het \omterm{def:g10:stats:mean}{gemiddelde} ($2820$) wordt door
het ene hoge loon omhooggetrokken en staat op niemands loonbrief.
\end{solution}

\begin{solution}{exo:g10:stats:8}
Totale lengte: $12 \times 168 + 18 \times 163 = 2016 + 2934 = 4950$ cm voor
$30$ leerlingen, dus is het \omterm{def:g10:stats:mean}{gemiddelde} $\frac{4950}{30} = 165$ cm. Het ligt
dichter bij het \omterm{def:g10:stats:mean}{gemiddelde} van de meisjes, omdat zij talrijker zijn (gewogen
\omterm{def:g10:stats:mean}{gemiddelde}, met gewichten $12$ en $18$).
\end{solution}

\begin{solution}{exo:g10:stats:9}
\emph{1.} Het nieuwe \omterm{def:g10:stats:mean}{gemiddelde} is
\[
\bar y = \frac{y_1 + \dots + y_N}{N}
= \frac{(a x_1 + b) + \dots + (a x_N + b)}{N}
= \frac{a(x_1 + \dots + x_N) + Nb}{N}
= a \bar x + b .
\]

\emph{2.} Met $a = 1.8$ en $b = 32$: \omterm{def:g10:stats:mean}{gemiddelde}
$1.8 \times 20 + 32 = 68$ graden Fahrenheit.
\end{solution}

\begin{solution}{pb:g10:stats:1}
\textbf{1.} \omterm{def:g10:stats:mean}{Gemiddelde}: $\frac{2\,550}{7} \approx 364$ duizend euro;
\omterm{def:g10:stats:median}{mediaan}: de vierde van de zeven gerangschikte prijzen, $230$. De \omterm{def:g10:stats:median}{mediaan}
beschrijft de straat --- zes van de zeven huizen kosten lang geen $364$.

\textbf{2.} Zonder het landhuis: \omterm{def:g10:stats:mean}{gemiddelde} $\frac{1\,350}{6} = 225$,
\omterm{def:g10:stats:median}{mediaan} $\frac{220 + 230}{2} = 225$. Het \omterm{def:g10:stats:mean}{gemiddelde} zakte met ongeveer
$139$; de \omterm{def:g10:stats:median}{mediaan} met $5$. Eén extreme waarde kan het \omterm{def:g10:stats:mean}{gemiddelde} overal heen
sleuren; de \omterm{def:g10:stats:median}{mediaan} merkt het nauwelijks --- ze is \emph{robuust}.

\textbf{3.} \omterm{def:g10:stats:mean}{Gemiddelde} $= \frac{0 \times 20 + 1 \times 25 + 2 \times 30 +
3 \times 15 + 4 \times 10}{100} = \frac{170}{100} = 1.7$. Geen enkel gezin
heeft $1.7$ kinderen, maar $1.7 \times 100 = 170$ \emph{is} het exacte
totale aantal kinderen: het \omterm{def:g10:stats:mean}{gemiddelde} legt het totaal vast, gelijk
herverdeeld --- zijn ware talent.

\textbf{4.} $N = 15$: \omterm{def:g10:stats:median}{mediaan} $=$ 8ste waarde $= 11$; $Q_1 =$ 4de waarde
$= 8$; $Q_3 =$ 12de waarde $= 14$.

\textbf{5.} \omterm{def:g10:stats:quartiles}{Variatiebreedte} $19 - 4 = 15$: de volledige spreiding, uitersten
inbegrepen. \omterm{def:g10:stats:quartiles}{Interkwartielafstand} $14 - 8 = 6$: de breedte van de middelste
helft --- spreiding zonder de ruis van de uitersten.

\textbf{6.} Een verschuiving met $+b$ verplaatst \emph{elke} waarde, dus
blijven verschillen van waarden --- \omterm{def:g10:stats:quartiles}{variatiebreedte} en \omterm{def:g10:stats:quartiles}{interkwartielafstand}
--- onveranderd: $(x_i + b) - (x_j + b) = x_i - x_j$. Een schaling met $a$
vermenigvuldigt alle verschillen met $a$: \omterm{def:g10:stats:quartiles}{variatiebreedte} en
\omterm{def:g10:stats:quartiles}{interkwartielafstand} worden met $\abs a$ geschaald.

\textbf{7.} Optellende aanpassing: $6 \to 9$ en $18 \to 21$ --- de zwakke
leerling wint $50\,\%$, de sterke $17\,\%$, en $21$ doet de schaal op $20$
maar net barsten. Vermenigvuldigende aanpassing: $6 \to 8$ en $18 \to 24$ ---
de verhoudingen blijven, maar $24$ kan niet op een schaal van $20$: het
plafondprobleem. Optellende aanpassingen vleien de onderkant;
vermenigvuldigende laten de bovenkant ontploffen.

\textbf{8.} Totaal aantal punten: $20 \times 12 + 30 \times 10 = 540$, dus is
het samengevoegde \omterm{def:g10:stats:mean}{gemiddelde} $\frac{540}{50} = 10.8$: de dertig leerlingen
van klas B wegen zwaarder dan de twintig van klas A --- \omterm{def:g10:stats:mean}{gemiddelden} voeg je
\emph{gewogen} samen, nooit met een gewoon \omterm{def:g10:stats:mean}{gemiddelde}.

\textbf{9.} \omterm{def:g10:stats:median}{Mediaan} van $A$: $2$; van $B$: $5$. Samengevoegde lijst
$\{0, 1, 2, 3, 10\}$: \omterm{def:g10:stats:median}{mediaan} $2$. Het naar omvang gewogen \omterm{def:g10:stats:mean}{gemiddelde} van de
medianen zou $\frac{3 \times 2 + 2 \times 5}{5} = 3.2 \neq 2$ zijn: medianen
dragen geen totaal, dus bestaat er geen formule om ze samen te voegen ---
je moet de hele lijst opnieuw sorteren.

\textbf{10.} \omterm{def:g10:stats:mean}{Gemiddelde}: $\frac{31.3}{5} = 6.26$. \omterm{def:g10:stats:median}{Mediaan}: $5.4$. Getrimd
\omterm{def:g10:stats:mean}{gemiddelde}: $\frac{5.3 + 5.4 + 5.5}{3} = 5.4$. Eén geestdriftig (of omgekocht)
jurylid verschoof het \omterm{def:g10:stats:mean}{gemiddelde} met bijna een punt; de \omterm{def:g10:stats:median}{mediaan} en het
getrimde \omterm{def:g10:stats:mean}{gemiddelde} haalden hun schouders op --- vandaar de trimregels in
beoordeelde sporten.

\textbf{11.} Fabriek A: middelste helft binnen $0.2$ mm; fabriek B: binnen
$4$ mm. Hetzelfde \omterm{def:g10:stats:mean}{gemiddelde}, tegengestelde betrouwbaarheid: koop bij A. De
\omterm{def:g10:stats:quartiles}{interkwartielafstand} besliste --- \omterm{def:g10:stats:mean}{gemiddelden} alleen vergelijken \omterm{prop:g10:coordgeom:midpoint}{middens},
nooit regelmaat.

\textbf{12.} Een \omterm{def:g10:stats:median}{mediaan} ver onder het \omterm{def:g10:stats:mean}{gemiddelde} verraadt een \emph{lange
rechterstaart}: de meeste wachttijden kort, enkele rampzalig. Voorbeeld:
$\{5, 5, 10, 10, 95\}$ --- \omterm{def:g10:stats:median}{mediaan} $10$, \omterm{def:g10:stats:mean}{gemiddelde} $\frac{125}{5} = 25$.

\textbf{13.} Het betekent dat $75\,\%$ van de baby's van die leeftijd korter
is (en $25\,\%$ langer): een rangorde, geen maat. ``Elk kind boven het
\omterm{def:g10:stats:mean}{gemiddelde}'' is onmogelijk --- overtroffen alle waarden het \omterm{def:g10:stats:mean}{gemiddelde}, dan
zou hun \omterm{def:g10:stats:mean}{gemiddelde} het \omterm{def:g10:stats:mean}{gemiddelde} overtreffen, en dat is het zelf:
tegenspraak. Maar \emph{bijna} alle kan: één heel klein kind kan het
\omterm{def:g10:stats:mean}{gemiddelde} onder alle andere houden. Boven de \emph{\omterm{def:g10:stats:median}{mediaan}}: nooit meer dan
de helft, per definitie. Rangordes laten zich niet vleien, \omterm{def:g10:stats:mean}{gemiddelden} wel.

\textbf{14.} Typisch punt: gelijk, \omterm{def:g10:stats:median}{mediaan} $11$ voor allebei. Homogeniteit
van het \omterm{prop:g10:coordgeom:midpoint}{midden}: B strakker (\omterm{def:g10:stats:quartiles}{interkwartielafstand} $3$ tegen $5$). Uitersten:
B wilder (\omterm{def:g10:stats:quartiles}{variatiebreedte} $17$ tegen $13$). B is een klas van gelijkaardige
leerlingen plus enkele uitschieters; A is gelijkmatig gespreid.

\textbf{15.} (1) Is $5\,000$ het \omterm{def:g10:stats:mean}{gemiddelde} of de \omterm{def:g10:stats:median}{mediaan} --- en mag ik de
\omterm{def:g10:stats:median}{mediaan}? (2) Wat is de spreiding (\omterm{def:g10:stats:quartiles}{kwartielen}), zodat ik weet hoe typisch
``typisch'' is? (3) Wie werd geteld --- steekproefgrootte en selectie (de les
over enquêtes uit de weekendopgave over hoe gegevens liegen in het
onderbouwvolume)?

\textbf{16.} Bijvoorbeeld $100, 150, 200, 250, 800$: \omterm{def:g10:stats:median}{mediaan} $200$,
\omterm{def:g10:stats:mean}{gemiddelde} $\frac{1\,500}{5} = 300$.

\textbf{17.} $\sum (x_i - \bar x) = \sum x_i - N \bar x = N\bar x - N\bar x
= 0$. Controle: $-200 - 150 - 100 - 50 + 500 = 0$. Het \omterm{def:g10:stats:mean}{gemiddelde} is het
punt waar de afwijkingen elkaar in evenwicht houden --- het zwaartepunt van
de gegevens (de $G$ uit de weekendopgave over het zwaartepunt in het
onderbouwvolume, in één dimensie).

\textbf{18.} Op de \omterm{def:g10:stats:median}{mediaan} $4$: $3 + 1 + 0 + 5 + 9 = 18$ km. Op het
\omterm{def:g10:stats:mean}{gemiddelde} $6$: $5 + 3 + 2 + 3 + 7 = 20$ km. Op $5$:
$4 + 2 + 1 + 4 + 8 = 19$ km. De \omterm{def:g10:stats:median}{mediaan} wint --- je ervan verwijderen
passeert altijd meer winkels dan het er nadert, zodat het totaal alleen maar
kan groeien.

\textbf{19.} $0.3^{10} \approx 6 \times 10^{-6}$: ongeveer zes piloten
\emph{per miljoen}. Bij $4\,000$ piloten is het verwachte aantal $0.02$ ---
nul, zoals luitenant Daniels vaststelde. De luchtmacht stopte met ontwerpen
voor de \omterm{def:g10:stats:mean}{gemiddelde} man en maakte alles regelbaar: stoelen, pedalen, riemen
--- ontwerp voor de \emph{spreiding}, niet voor het \omterm{prop:g10:coordgeom:midpoint}{midden}.

\textbf{20.} Totalen: het \emph{\omterm{def:g10:stats:mean}{gemiddelde}}. Typisch individu: de
\emph{\omterm{def:g10:stats:median}{mediaan}}. Eerlijkheid van de spreiding: de \emph{\omterm{def:g10:stats:quartiles}{kwartielen}} (en de
\omterm{def:g10:stats:quartiles}{interkwartielafstand}). Besmetting: het \emph{getrimde \omterm{def:g10:stats:mean}{gemiddelde}}. En de
slotregel: de \omterm{def:g10:stats:mean}{gemiddelde} mens bestaat niet --- maar de mediane mens bijna
wel.
\end{solution}
